% Minimal Preamble - Version 2.1
% Technical book template optimized for code and diagrams
% Last updated: 2026-01-09

% ========== FONTS AND ENCODING ==========
\usepackage[T1]{fontenc}
\usepackage[utf8]{inputenc}
\usepackage{lmodern}

% ========== PAGE LAYOUT ==========
\usepackage[margin=1in]{geometry}

% ========== COLORS (before other packages that use colors) ==========
\usepackage[dvipsnames]{xcolor}
\definecolor{rustblue}{RGB}{0, 119, 255}
\definecolor{codebg}{RGB}{248, 248, 248}

% ========== CODE LISTINGS ==========
% minted provides fancyvrb, so we get it for free
\usepackage{minted}
\usemintedstyle{friendly}
\setminted{
    fontsize=\footnotesize,
    baselinestretch=0.9,
    bgcolor=codebg,
    linenos,
    numbersep=5pt,
    frame=none,
    framesep=10pt,
    breaklines=true,
    breakanywhere=true
}

% ========== HEADERS AND FOOTERS ==========
\usepackage{fancyhdr}
\pagestyle{fancy}
\fancyhf{}
\fancyhead[LE]{\thepage\quad\leftmark}
\fancyhead[RO]{\rightmark\quad\thepage}
\fancyfoot[C]{}
\renewcommand{\headrulewidth}{0.4pt}
\setlength{\headheight}{25.5pt}  % Fix fancyhdr warning

% ========== GRAPHICS ==========
\usepackage{graphicx}
\usepackage{tikz}
\usetikzlibrary{arrows.meta}    % For arrow tips like Latex
\usetikzlibrary{positioning}    % For positioning syntax (of=)
\usepackage{forest}      % For directory tree diagrams

% ========== TABLES ==========
\usepackage{array}
\usepackage{booktabs}    % Professional tables

% ========== MARGIN NOTES ==========
\usepackage{marginnote}
\reversemarginpar
\setlength{\marginparwidth}{1.2in}

% ========== BIBLIOGRAPHY ==========
\usepackage[style=ieee, backend=biber]{biblatex}

% ========== GLOSSARY ==========
\usepackage[toc, nopostdot]{glossaries}

% ========== TABLE OF CONTENTS ==========
\usepackage{minitoc}
\dominitoc

% ========== TITLE FORMATTING ==========
\usepackage{titlesec}
% Chapter titles - slightly smaller
\titleformat{\chapter}[block]
  {\Large\bfseries}
  {\thechapter.\quad}
  {0pt}
  {\LARGE}

% Section titles - smaller
\titleformat{\section}
  {\large\bfseries}
  {\thesection\quad}
  {0pt}
  {}

% Subsection titles - smaller
\titleformat{\subsection}
  {\normalsize\bfseries}
  {\thesubsection\quad}
  {0pt}
  {}

% ========== HYPERLINKS (load late to avoid conflicts) ==========
\usepackage[pdfborder={0 0 0}, breaklinks=true]{hyperref}
\hypersetup{
    colorlinks=true,
    linkcolor=rustblue,
    filecolor=magenta,
    urlcolor=rustblue,
    pdftitle={Implementing Rust Support for Embedded Hardware},
    pdfauthor={Marcelo Correa},
    pdfsubject={Embedded Systems, Rust Programming},
    pdfkeywords={Rust, Embedded, EFR32MG24, HAL, PAC, Wireless}
}

% ========== CUSTOM ENVIRONMENTS ==========
\usepackage{mdframed}
\usepackage{enumitem}

\newenvironment{example}
  {\begin{mdframed}[backgroundcolor=gray!10, linewidth=0.5pt]}
  {\end{mdframed}}

\newenvironment{notebox}
  {\begin{mdframed}[backgroundcolor=blue!10, linewidth=0.5pt]}
  {\end{mdframed}}

\newenvironment{warningbox}
  {\begin{mdframed}[backgroundcolor=red!10, linewidth=0.5pt]}
  {\end{mdframed}}

\newenvironment{tipbox}
  {\begin{mdframed}[backgroundcolor=green!10, linewidth=0.5pt]}
  {\end{mdframed}}

\newenvironment{remark}
  {\begin{mdframed}[backgroundcolor=yellow!10, linewidth=0.5pt]}
  {\end{mdframed}}

\newenvironment{keytakeaways}
  {\begin{itemize}[leftmargin=*]}
  {\end{itemize}}

\newenvironment{exercises}
  {\begin{enumerate}[leftmargin=*]}
  {\end{enumerate}}

\newenvironment{exercise}[1][Easy]
  {\paragraph{Exercise (#1):}}
  {}

% ========== CUSTOM COMMANDS ==========
% Inline code with Rust syntax highlighting
\newcommand{\rustcode}[1]{\mintinline{rust}{#1}}

% Code inline - uses fancyvrb's \Verb for proper verbatim handling
% This handles underscores, hashes, and all special characters
\usepackage{xparse}
\NewDocumentCommand{\codeinline}{v}{%
  \texttt{#1}%
}

% Crate names (same verbatim handling)
\NewDocumentCommand{\crate}{v}{%
  \texttt{#1}%
}

% File paths (same verbatim handling)
\NewDocumentCommand{\filepath}{v}{%
  \texttt{#1}%
}

% Structure/concept names (just italic, no special handling)
\newcommand{\structure}[1]{\textit{#1}}

% Chapter thumb placeholder
\newcommand{\chapthumb}[1]{}

% Listings list stub (we use minted, not listings package)
\newcommand{\lstlistoflistings}{%
  \chapter*{List of Code Listings}%
  \markboth{List of Code Listings}{List of Code Listings}%
}

% ========== NUMBERING DEPTH ==========
\setcounter{secnumdepth}{2}  % Number sections and subsections
\setcounter{tocdepth}{2}     % Show sections and subsections in TOC

% ========== UNICODE CHARACTERS THE BOOK USES ==========
% T1/lmodern does not carry every glyph the text uses: an em-dash or a bullet
% disappears silently (pdflatex also refuses them outright). Mapping them to LaTeX
% commands is one fix in one place, and it keeps the source readable.
\usepackage{newunicodechar}
\newunicodechar{—}{---}
\newunicodechar{–}{--}
\newunicodechar{•}{\textbullet}
\newunicodechar{…}{\ldots}
\newunicodechar{×}{\ensuremath{\times}}
\newunicodechar{→}{\ensuremath{\rightarrow}}
\newunicodechar{≈}{\ensuremath{\approx}}

% ========== MONOSPACE FONT WITH THE GLYPHS THE EXAMPLES USE ==========
% The directory trees in the code blocks use box-drawing characters (U+2500..2514)
% that T1/lmodern does not carry: they vanish without an error, leaving a tree of
% bare names. DejaVu Sans Mono has them, and XeLaTeX can select it directly.
\usepackage{fontspec}
\setmonofont{DejaVu Sans Mono}[Scale=0.86]
